\section{Proposed Methodology}

We will design a traditional RTS game interface and an LLM interface for controlling a MAS in a game environment.  We will then perform a \textbf{within-subjects} user study to evaluate the performance and user experience of the two interfaces, with the ordering of the interfaces \textbf{counterbalanced} for each participant.  A within-subjects design serves the function of controlling for the \textbf{confounding variable} of RTS game experience.  Moreover, by counterbalancing the order of the interfaces' presentation to the participant, we seek to abate any \textbf{learning} or \textbf{sequence effects}.

We will evaluate a participant's performance with the MAS interface in terms of the score they achieve in the game environment.  We will evaluate a participant's user-experience with the MAS interface in terms of task load and usability.  To measure task load, we will be using the \textbf{NASA Task Load Index (NASA-TLX)} \cite{hart_development_1988}.  To measure usability, we will be using the \textbf{System Usability Scale (SUS)} \cite{brooke_susa_1996}.  Additionally, we will be collecting qualitative feedback from the participants in the form of a \textbf{Semi-Structured Interview} to gain insight into their experience with the two interfaces.  Our \textbf{research hypothesis} is that the LLM interface will lead to a decreased task load and no meaningful difference in performance when compared to the RTS game interface.

The first major part of our project is creating a game environment for the MAS to operate in.  We will be using the Robotarium \cite{pickem_robotarium_2017} as our MAS platform and we will utilize the Robotarium's overhead projector to create a game environment for the MAS to operate within.  While we have not yet decided specific elements of the game environment, it should satisfy the following:\begin{itemize}[noitemsep, topsep=0pt]
    \item challenging enough to induce task saturation at times
    \item require the simultanous coordination of a large number of robots
    \item include both positive rewards to serve as goals and negative rewards to serve as constraints
\end{itemize}
Much of the previous work has focused on contrived examples that do not require conforming to a set of environmental constraints, whereas these are incredibly important aspects of real-world MAS.  We will be creating a scoring system for the game that will provide positive feedback to the MAS for completing tasks and negative feedback for failing to complete tasks or violating constraints.

With the Robotarium providing a software simulator, hardware, and runtime environment for the MAS, we will be able to focus on creating the game environment and interfaces for controlling the MAS.  The team has not yet decided on the specific LLM for use in the project, but selecting the LLM will be one of the first tasks for the project.  Additionally, we will be creating a traditional RTS game interface for controlling the MAS.  A technology stack has also not yet been selected for the RTS game interface, but we intend for the design to adhere to the human-computer interaction principles for RTS games as presented in \cite{wenjun_hci_2006} and \cite{riaz_user_2021}.
